\begin{lemma}[4.8.8]
\label{IV.2.4.8.8-fr}
Con las notaciones de \textup{(4.8.2)}:
\begin{enumerate}
\item[\textup{(i)}] Sea
$\overline u:\mathscr F\otimes_kK\to\mathscr G\otimes_kK$ un
$\mathcal O_{X_{(K)}}$-homomorfismo, y sea
$\overline{\mathscr H}\subset
(\mathscr F\otimes_kK)\oplus(\mathscr G\otimes_kK)$ su grafo. Para que
$\overline u$ esté definido sobre $K'$, es necesario y suficiente que
$\overline{\mathscr H}$ lo esté.
\item[\textup{(ii)}] Sea $\overline Z$ un subpreesquema cerrado de $X_{(K)}$, y sea
$\overline{\mathscr J}$ el ideal cuasicoherente de $\mathcal O_{X_{(K)}}$
que define $\overline Z$. Para que $\overline Z$ esté definido sobre $K'$, es
necesario y suficiente que $\overline{\mathscr J}$ lo esté.
\item[\textup{(iii)}] Sea $\overline f$ un $K$-morfismo:
$X_{(K)}\to Y_{(K)}$ y sea $\overline Z$ el subpreesquema de
$X_{(K)}\times_KY_{(K)}$, grafo de $\overline f$ \textup{(\textbf I, 5.3.11)}; para que
$\overline f$ esté definido sobre $K'$, es necesario y suficiente que $\overline Z$ lo
esté.
\end{enumerate}
\end{lemma}

\begin{proof}
La afirmación (ii) resulta inmediatamente de (\textbf I, 4.4.5). La necesidad en la
afirmación (iii) es evidente; supongamos recíprocamente que
$\overline Z=Z'_{(K)}$, donde $Z'$ es un subpreesquema de
$(X\times_kY)_{(K')}$; si $p':Z'\to X_{(K')}$ es la restricción a $Z'$ de la
primera proyección, se sabe que $p'_{(K)}$ es un isomorfismo de preesquemas,
por tanto lo mismo ocurre con $p'$ (2.7.1, (viii)); pero esto significa (\textbf I, 5.3.11)
que $Z'$ es el grafo de un $K'$-morfismo
$f':X_{(K')}\to Y_{(K')}$, y se tiene entonces
$\overline f=f'_{(K)}$ (\textbf I, 5.3.12).

Finalmente, para demostrar (i), se puede, en virtud de (4.8.5), suponer que
$X=\operatorname{Spec}(A)$ es afín, $\mathscr F=\widetilde M$,
$\mathscr G=\widetilde N$, donde $M$ y $N$ son $A$-módulos, y
$\overline u=(\overline v)^{\sim}$, donde
$\overline v:M_{(K)}\to N_{(K)}$ es un $A_{(K)}$-homomorfismo. Supongamos
que el grafo de $\overline v$ sea de la forma $P'_{(K)}$, donde $P'$ es un
sub-$A_{(K')}$-módulo de $M_{(K')}\oplus N_{(K')}$; si
$p':P'\to M_{(K')}$ es la restricción a $P'$ de la primera proyección,
se sabe que $p'\otimes1_K$ es un isomorfismo de $A_{(K)}$-módulos, por tanto,
por planitud fiel $(\mathbf{0}_{\mathrm I},6.4.1)$, $p'$ es un isomorfismo de
$A_{(K')}$-módulos, lo que prueba que $P'$ es el grafo de un
$A_{(K')}$-homomorfismo $v':M_{(K')}\to N_{(K')}$ y se tiene evidentemente
$\overline v=v'\otimes1_K$.
\end{proof}

\begin{proposition}[4.8.9]
\label{IV.2.4.8.9-fr}
Sean $k$ un cuerpo, $K$ una extensión de $k$, $X$ un $k$-preesquema,
$\mathscr F$ un $\mathcal O_X$-módulo cuasicoherente. Sea $\overline{\mathscr H}$ un
$\mathcal O_{X_{(K)}}$-submódulo cuasicoherente de
$\mathscr F\otimes_kK$. Entonces existe un menor cuerpo de definición de
$\overline{\mathscr H}$.
\end{proposition}

\begin{proof}
Cuando $X$ es afín, la afirmación no es sino (4.8.7). En el caso general,
consideremos un recubrimiento $(X_\alpha)$ de $X$ por abiertos afines; en
virtud de lo anterior, existe para todo $\alpha$ un menor subcuerpo $K'_\alpha$
de $K$ que contiene $k$ y tal que
$\overline{\mathscr H}|(X_\alpha)_{(K)}$ esté definido sobre $K'_\alpha$. En virtud de
(4.8.5), el subcuerpo de $K$ engendrado por los $K'_\alpha$ es el menor
cuerpo de definición de $\overline{\mathscr H}$.
\end{proof}

\begin{corollary}[4.8.10]
\label{IV.2.4.8.10-fr}
Con las notaciones de \textup{(4.8.9)}, sea $\mathscr G$ un segundo
$\mathcal O_X$-módulo cuasicoherente, y sea
$\overline u:\mathscr F\otimes_kK\to\mathscr G\otimes_kK$ un
$\mathcal O_{X_{(K)}}$-homomorfismo; entonces existe un menor cuerpo de
definición de $\overline u$.
\end{corollary}

\oldpage[IV-2]{83}

\begin{proof}
En efecto, resulta de (4.8.8) que tal cuerpo es también un menor cuerpo de
definición del grafo de $\overline u$, y la existencia de este último resulta
de (4.8.9).
\end{proof}
